\documentclass[a4paper]{article}

\usepackage{graphicx}
\usepackage{amsmath,amssymb}
\usepackage{minted}
\usepackage{multirow}
\usepackage{float}
\usepackage{algorithm}
\usepackage{algpseudocode}
\usepackage{todonotes}
\usepackage{forest}
\usepackage{xstring}
\usepackage{xcolor}
\usepackage[most]{tcolorbox}
\usepackage{xparse}
\usepackage{natbib}
\usepackage{adjustbox}
\usepackage[active,tightpage]{preview}

\PreviewBorder=4pt

\title{Source formula/table live previews and Visual display-math sizing}
\date{2026-09-17}
\author{Mohammad Rahmani}

\begin{document}
\maketitle

\section{Fix}
The Visual display-math size policy is now applied to the rendered image itself
before the image reaches the QTextDocument projection.  The previous code
computed readable logical dimensions but the display insertion path never used
those dimensions, so a standalone equation could still be shown at the tiny
native raster size.  One-line, tall, and multi-row display formulas now receive
font-relative readable dimensions while remaining bounded by the Visual
viewport width.

\section{Source live previews}
Source mode gains non-destructive graphical projections for inline mathematics,
display mathematics, and tabular environments.  The canonical LaTeX text is
never modified by a preview: overlays are anchored to the source range and
space is reserved only in QTextBlock formatting.  Inline formula renders are
anchored immediately after the formula source position; display formulas and
tables are centered below the source range.

The Source controller listens to document/content/view changes and uses a short
debounce.  LuaLaTeX renders only uncached snippet text, so editing a formula or
table causes the graphical projection to refresh after the source changes while
unchanged snippets reuse their in-memory images.  A failing unsupported snippet
is isolated so it does not suppress valid previews elsewhere in the document.

\end{document}
